\documentclass[11pt]{article}
% Shared preamble for Theseus user guides.
% Compiled with pdflatex. See docs/guides/BUILD.md for rationale.

\usepackage[a4paper,margin=0.85in]{geometry}
\usepackage[T1]{fontenc}
\usepackage[utf8]{inputenc}

% Body: Palatino (via mathpazo). Headings: Helvetica (sans).
% Palatino is requested in the house style as a substitute for Charter
% because Charter is not part of the standard TeX Live distribution
% that the build container ships.
\usepackage{mathpazo}
\usepackage[scaled=0.92]{helvet}
\usepackage{courier}
\renewcommand{\familydefault}{\rmdefault}

\usepackage{microtype}
\usepackage{parskip}
\usepackage{enumitem}
\usepackage{booktabs}
\usepackage{longtable}
\usepackage{titlesec}
\usepackage{xcolor}
\usepackage{fancyhdr}
\usepackage{fancyvrb}
\usepackage{graphicx}
% Allow `_` in text mode without escaping. The screenshot file names and
% some environment-variable references in the guides contain underscores.
\usepackage{underscore}
\usepackage{tcolorbox}
\tcbuselibrary{breakable,skins}
\usepackage[hidelinks,colorlinks=true,linkcolor=black,urlcolor=teal,citecolor=black]{hyperref}

% Sans for chapter and section headings; serif body keeps prose readable
% on screen and in print.
\titleformat{\section}{\sffamily\Large\bfseries}{\thesection}{0.7em}{}
\titleformat{\subsection}{\sffamily\large\bfseries}{\thesubsection}{0.6em}{}
\titleformat{\subsubsection}{\sffamily\normalsize\bfseries}{\thesubsubsection}{0.5em}{}

\pagestyle{fancy}
\fancyhf{}
\fancyhead[L]{\small\sffamily Theseus User Guide}
\fancyhead[R]{\small\sffamily\thepage}
\renewcommand{\headrulewidth}{0.3pt}

\setlength{\parindent}{0pt}

% Convenience commands shared by every guide.
\newcommand{\page}[1]{\textbf{#1}}
\newcommand{\file}[1]{\texttt{#1}}
\newcommand{\route}[1]{\texttt{#1}}
\newcommand{\env}[1]{\texttt{#1}}
\newcommand{\code}[1]{\texttt{#1}}

% Callout boxes used by every guide.
\newtcolorbox{tldr}{
  colback=teal!5!white,colframe=teal!55!black,
  fonttitle=\sffamily\bfseries,title=TL;DR,
  breakable,boxrule=0.4pt,arc=2pt
}
\newtcolorbox{youwillneed}{
  colback=gray!4!white,colframe=gray!55!black,
  fonttitle=\sffamily\bfseries,title=You will need,
  breakable,boxrule=0.4pt,arc=2pt
}
\newtcolorbox{notcovered}{
  colback=orange!5!white,colframe=orange!55!black,
  fonttitle=\sffamily\bfseries,title=This guide does NOT cover,
  breakable,boxrule=0.4pt,arc=2pt
}
\newtcolorbox{preview}{
  colback=yellow!8!white,colframe=yellow!55!black,
  fonttitle=\sffamily\bfseries,title=Preview --- partially shipped,
  breakable,boxrule=0.4pt,arc=2pt
}
\newtcolorbox{nextup}{
  colback=blue!4!white,colframe=blue!55!black,
  fonttitle=\sffamily\bfseries,title=Where to read next,
  breakable,boxrule=0.4pt,arc=2pt
}

% Insert a screenshot only when the file exists, so the build does not
% break before the Playwright capture run is wired up.
\newcommand{\screenshot}[3]{%
  \IfFileExists{screenshots/#1}{%
    \begin{figure}[h]\centering
      \includegraphics[width=\linewidth]{screenshots/#1}
      \caption{#2}\label{fig:#3}
    \end{figure}%
  }{%
    \begin{center}\small\itshape%
      [screenshot \texttt{screenshots/#1} not yet captured --- #2]%
    \end{center}%
  }%
}


\title{\sffamily\bfseries Knowledge and Principles\\[2pt]\large Guide 2 of 6}
\date{}
\author{}

\begin{document}
\maketitle
\thispagestyle{empty}

\begin{tldr}
Theseus turns uploaded material into firm principles through a four-stage
pipeline: ingest, claim extraction, clustering, and principle
distillation. The first two stages run automatically; the last two
involve a triage queue you sit in front of when you want to keep the
firm's working memory clean. This guide walks the pipeline end to end,
shows where each artifact lives, and explains the founder's role in
triage.
\end{tldr}

\begin{youwillneed}
\begin{itemize}[leftmargin=*,itemsep=0.2em,topsep=0.2em]
  \item A signed-in Codex session.
  \item At least one upload in the corpus --- if you do not have one,
        complete Guide 1 first.
  \item Optional: a couple of related documents already uploaded, so
        that clustering has something to merge across.
\end{itemize}
\end{youwillneed}

\begin{notcovered}
\begin{itemize}[leftmargin=*,itemsep=0.2em,topsep=0.2em]
  \item Question-answering with the Oracle --- Guide 3.
  \item Currents opinions that cite principles --- Guide 4.
  \item Forecasts that may use principle-backed conclusions as a
        reasoning step --- Guide 5.
  \item Operator-only deletion / retraction of principles --- Guide 6.
\end{itemize}
\end{notcovered}

\section{The pipeline in one diagram}

\begin{center}
\begin{tabular}{l@{\;\(\to\)\;}l@{\;\(\to\)\;}l@{\;\(\to\)\;}l}
\sffamily Upload & \sffamily Claims & \sffamily Clusters & \sffamily Principles \\
ingest layer & claim extractor & clustering pass & distillation
\end{tabular}
\end{center}

Each arrow above corresponds to a worker stage. Stages 1 and 2 run on
upload. Stages 3 and 4 run on a schedule and are also triggerable from
the founder triage queue.

\section{Stage 1 --- ingest}

Uploads enter through \page{Upload} on the Codex
(\route{/(authed)/upload}). The ingest worker normalizes the file
(MIME detection, encoding repair, page extraction for PDFs), persists
it to storage, and writes an \texttt{Upload} row with
\texttt{publishedAt=NULL}, \texttt{visibility='private'}, and a
processing badge that the UI watches.

What you see in the Codex: the new row appears in
\page{Upload} list, then in \page{Knowledge} once at least one claim
has been extracted from it. Provenance is captured at ingest:
file hash, source URL (if any), uploader identity, and timestamps.

\section{Stage 2 --- claim extraction}

The claim extractor reads the normalized upload and emits short,
declarative, falsifiable statements. Each claim carries:

\begin{itemize}[leftmargin=*,itemsep=0.2em]
  \item the upload it came from (provenance);
  \item a topic label and an embedding vector for clustering;
  \item the verbatim span it summarizes (so the Oracle can cite it);
  \item an extractor-confidence score.
\end{itemize}

Look at any claim in \page{Knowledge} and you can click straight back
to the source line. This is the contract that makes everything
downstream auditable: no claim is allowed to exist without a pointer
to the text that produced it.

\screenshot{02_knowledge_claim_view.png}{The claim view in \page{Knowledge}, with a single claim expanded to show its source span.}{kn-claim}

\section{Stage 3 --- clustering}

The clustering pass groups claims by topic and semantic similarity.
The output is a candidate cluster: a small set of claims that look
like they are talking about the same thing.

Clusters are not yet principles. They are an intermediate object: an
invitation to consider whether the firm wants to lift the underlying
material into a reusable principle.

\section{Stage 4 --- principle distillation}

For each candidate cluster the distillation worker produces a draft
\page{Principle}: a single, sharp statement that the underlying claims
support, plus an explicit list of the supporting claims and an
explicit list of the strongest objections found in the corpus.

A principle has three states that matter:

\begin{itemize}[leftmargin=*,itemsep=0.2em]
  \item \texttt{draft} --- not visible publicly. Visible in the triage
        queue.
  \item \texttt{accepted} --- vetted. May or may not be public.
  \item \texttt{rejected} / \texttt{merged} / \texttt{needs\_rereview}
        --- terminal or paused states.
\end{itemize}

A principle becomes publicly visible only when
\texttt{status = 'accepted'} \textbf{and} \texttt{publicVisible = true}.
Both are required. See \file{docs/operator/public\_surfacing.md} for
the authoritative predicate table.

\section{The founder's role --- the triage queue}

The principles surface in the Codex
(\route{/(authed)/principles}) is the triage queue. Open it and you
will see candidate principles in three columns:

\begin{itemize}[leftmargin=*,itemsep=0.2em]
  \item \textbf{New drafts.} Just produced by distillation. Read the
        principle, scan supporting claims, scan objections. Decide:
        \texttt{accept}, \texttt{merge into existing}, or
        \texttt{reject with reason}.
  \item \textbf{Needs re-review.} Either contradicted by a newer
        cluster or flagged by the adversarial pass. The previous
        decision should be re-justified or rolled back.
  \item \textbf{Accepted.} The firm's working memory. Editable in
        place; toggling \texttt{publicVisible} sends a principle to
        (or removes it from) the public methodology page.
\end{itemize}

\screenshot{02_principles_triage.png}{The triage queue, with the three columns described above.}{pr-triage}

The triage queue is the only stage in this pipeline that requires
human judgment. If you ignore it for a week, drafts accumulate; the
firm's automatic answers keep working off the last batch of accepted
principles. Nothing collapses, but the firm's public surface stops
evolving.

\section{Where principles live afterwards}

\begin{itemize}[leftmargin=*,itemsep=0.2em]
  \item \textbf{Codex (private).} \route{/(authed)/principles}. The
        full draft + accepted + rejected set, with editing controls.
  \item \textbf{Public methodology page.} \route{/methodology/principles}
        on the public site. The accepted-and-public subset.
  \item \textbf{Currents opinions.} An opinion may cite principles
        when it leans on a firm position rather than a one-off claim.
  \item \textbf{Oracle answers.} The Oracle may cite principles for
        the same reason.
\end{itemize}

\section{Common failure modes}

\begin{itemize}[leftmargin=*,itemsep=0.3em]
  \item \textbf{Drafts pile up.} The clustering pass keeps proposing
        principles; nobody reviews them. Reserve a recurring slot
        weekly for triage. Five to ten minutes per draft is realistic.
  \item \textbf{A principle on the public site contradicts a private
        accepted one.} Flip \texttt{publicVisible=false} on the public
        one and re-triage. Both stay in the corpus; only their public
        surfacing changes.
  \item \textbf{The cluster is junk.} Reject the draft with a reason.
        The reason is read by the next distillation pass, which uses
        it to avoid re-proposing the same cluster.
\end{itemize}

\section{Where the docs are}

\begin{itemize}[leftmargin=*,itemsep=0.2em]
  \item \file{docs/methods/MQS\_Specification.md} --- the formal claim
        / conclusion / principle specification.
  \item \file{docs/operator/public\_surfacing.md} --- which fields gate
        each public surface, including principles.
  \item \file{noosphere/noosphere/distillation/} --- the worker code,
        for the curious.
\end{itemize}

\begin{nextup}
Continue with \href{03_The_Oracle.pdf}{Guide 3 --- The Oracle} for how
question-answering uses the corpus you have just helped curate.
\end{nextup}

\end{document}
